\SourceLeaf{374}%
from \(AO\) a segment \(OM=NL\) and draws \(ML\), and if at \(M\) one again constructs on \(MA\) an angle equal to angle \(A\) whose other arm \(MC\) meets \(AN\) at the point \(C\) between \(A\) and \(L\), then cuts off from \(AM\) a segment \(MB=LC\) and draws \(BC\), and continues this construction in the same way so that on the line \(OA\) lie the points \(O,M,B,E,G,K\) and so on, and on the line \(NA\) lie the points \(N,L,C,D,F,H\) and so on, then it is asserted that the segments \(OM,MB,BE,EG,GK\) and so on, or the segments equal to them respectively, \(NL,LC,CD,DF,FH\) and so on, form a decreasing progression.

The author tries to prove this theorem by contradiction. He lists the other possible cases if the theorem were false and attempts to show that each is impossible. He claims that on that assumption one of the following five cases must occur. The successive segments, counting from \(OM\), would be
\par\noindent\hspace*{\parindent}1) all equal to one another, or
\par\noindent\hspace*{\parindent}2) each succeeding segment larger than the preceding one, or
\par\noindent\hspace*{\parindent}3) some equal to one another and the next larger or smaller, or
\par\noindent\hspace*{\parindent}4) some consecutive segments decreasing progressively and the ones that follow increasing progressively or
\par\noindent\hspace*{\parindent}5) the segments alternately larger and smaller.
\par\noindent This list leaves out the possible case in which the segments first increase progressively and then decrease progressively. In the reviewer's own conviction (the deeper reasons for which cannot be given here), disposing of that case would be the main issue and the real untying of the Gordian knot. Still, one can grant that the omission matters little here insofar as the author's way of proving the third case impossible, if it were valid, would extend of itself to this case too. Yet we cannot grant any validity to this very purported proof that the third case is impossible. The author puts it as follows. If, for example, in the third case the first two segments are assumed equal and the third larger, then \(DC\) would be greater than \(CL\). But \(AML\) is likewise an isosceles triangle that meets the same basic condition as the original triangle \(AON\). So if that third case with its assumed subdivision were the applicable one,
